\documentclass[12pt,a4paper,oneside]{article}
\usepackage[brazil]{babel}
\usepackage[utf8]{inputenc}
\usepackage{olymp}
\usepackage{color}
\usepackage[dvipsnames]{xcolor}
\usepackage{listings}

\setcounter{problem}{2}

\begin{document}

\begin{problem}{Produto escalar}{prodesc}{2}

Em álgebra linear, o produto escalar de dois vetores $\mathbf{A}$ e $\mathbf{B}$ (denotado por $\mathbf{A} \cdot \mathbf{B}$ ) é uma função que fornece um número real como resultado. Matematicamente, o produto escalar é definido como uma função $\mathbb{R}^n \times \mathbb{R}^n \rightarrow \mathbb{R}$ da seguinte forma:

\begin{equation}
    \mathbf{A} \cdot \mathbf{B} = \sum_{i=1}^{n} a_ib_i \notag
\end{equation}

sendo $N$ o números de elementos dos vetores e $a_i$ e $b_i$ seus elementos.

Por exemplo, se $\mathbf{A} = [1,2,3,4]$ e $\mathbf{B} = [5,6,7,8]$, temos que $n=4$  e que:

\begin{equation}
    \mathbf{A} \cdot \mathbf{B} = (1 \times 5) + (2 \times 6) + (3 \times 7) + (4 \times 8) = 70 \notag
\end{equation}

Neste problema você deverá fazer um programa para calcular o produto escalar entre dois vetores de números inteiros. Assim, o resultado obtido também será sempre um número inteiro.

%\Notes
%
%Observações, se tiver.

\InputFile

A entrada contém três linhas: a primeira contém um número inteiro $N$, indicando a quantidade de elementos dos vetores; a segunda contém os $N$ elementos do vetor $\mathbf{A}$; e a terceira os $N$ elementos do vetor $\mathbf{B}$.

Restrições: $1 \leq N \leq 1000$, $-100 \leq a_i, b_i \leq 100$.



\OutputFile

Seu programa deve gerar apenas uma linha de saída, contendo o resultado do produto escalar $\mathbf{A} \cdot \mathbf{B}$ .

% \Notes
% Escreva espaço em branco para separar os números, mas não escreva espaço em branco antes do primeiro nem depois do último número da lista. Após o último escreva apenas um caractere de quebra de linha.


\Example

\begin{example}
\exmp{
4
1 2 3 4
5 6 7 8
}{
70
}%
\end{example}

\begin{example}
\exmp{
3
1 3 -5
4 -2 1
}{
-7
}%
\end{example}

\begin{example}
\exmp{
9
0 1 2 3 4 5 6 7 8
0 1 2 3 4 5 6 7 8
}{
204
}%
\end{example}



%Comentários sobre o exemplo, se necessário.

\end{problem}

\end{document}
